\question
\subsection*{پردازه‌ها در لینوکس}

در این سؤال، با اجرای یک دستور ساده، شناسهٔ پردازه، رابطهٔ پدر و فرزند
و وضعیت پردازه را در لینوکس مشاهده می‌کنیم. برای انجام تمرین، یک محیط
لینوکس روی سیستم اصلی یا \lr{Virtual Machine} و ابزار
\lr{\texttt{ps}} کافی است؛ نیازی به نوشتن برنامهٔ \lr{C} ندارید.

\paragraph{بخش اول: اجرای پردازه و مشاهدهٔ مشخصات آن}
در یک پایانهٔ \lr{Bash}، دستورهای زیر را به‌ترتیب اجرا کنید. تمام
مراحل را در همان پایانه انجام دهید:

\begin{latin}
\begin{verbatim}
sleep 120 &
child=$!
echo "Shell PID: $$"
echo "Child PID: $child"
ps -o pid,ppid,stat,comm -p "$$,$child"
\end{verbatim}
\end{latin}

دستور \lr{\texttt{sleep}} برای مدت مشخصی منتظر می‌ماند و سپس تمام
می‌شود. قرار دادن \lr{\texttt{\&}} در انتهای دستور، آن را در پس‌زمینه
اجرا می‌کند تا بتوانید پیش از پایان آن، دستورهای بعدی را وارد کنید.
متغیر \lr{\texttt{child}} شناسهٔ همین پردازه را نگه می‌دارد.

\paragraph{بخش دوم: بررسی اطلاعات پردازه}
پیش از پایان زمان انتظار، با دستور زیر اطلاعات همان فرزند را از
\lr{\texttt{/proc}} بخوانید:

\begin{latin}
\begin{verbatim}
grep -E '^(Name|State|Pid|PPid):' "/proc/$child/status"
\end{verbatim}
\end{latin}

مقادیر شناسهٔ پردازه و شناسهٔ پدر را با خروجی
\lr{\texttt{ps}} مقایسه کنید. سپس برای پایان دادن به همین آزمایش،
دستورهای زیر را اجرا کنید:

\begin{latin}
\begin{verbatim}
kill "$child"
wait "$child"
\end{verbatim}
\end{latin}

دستور \lr{\texttt{kill}} در اینجا به پردازهٔ ساخته‌شدهٔ خودتان درخواست
پایان می‌دهد و \lr{\texttt{wait}} وضعیت پایان آن را دریافت می‌کند.
ممکن است پوسته پیام پایان پردازه را نمایش دهد و
\lr{\texttt{wait}} مقدار غیرصفر برگرداند؛ این نتیجه در این آزمایش طبیعی است.
اگر پردازه پیش از تکمیل مشاهده‌ها به‌طور عادی تمام شد، آزمایش را از
ابتدا تکرار کنید؛ شناسهٔ اجرای قبلی را دوباره استفاده نکنید.

\paragraph{پرسش‌های گزارش}
خروجی دستورها را ثبت کنید و به دو پرسش زیر کوتاه پاسخ دهید:

\begin{enumerate}
  \item تفاوت \lr{\texttt{PID}} و \lr{\texttt{PPID}} چیست؟ با استفاده
  از خروجی خود مشخص کنید پدر پردازهٔ \lr{\texttt{sleep}} کدام است و
  آیا دو روش مشاهده، شناسه‌های یکسانی نشان می‌دهند؟
  \item وضعیت مشاهده‌شده برای \lr{\texttt{sleep}} چیست؟ توضیح دهید
  این پردازه منتظر چه چیزی است و این وضعیت به کدام حالت از مدل
  پردازه در درس مربوط می‌شود.
\end{enumerate}

\paragraph{ویدیوی پاسخ و توضیح دستورها}
یک ویدیو با تصویر خوانای صفحه ضبط کنید؛ زمان پیشنهادی آن حدود دو تا
سه دقیقه است. اجرای آزمایش و خروجی‌های خود را نشان دهید و کاربرد
\lr{\texttt{\$\$}} و \lr{\texttt{\$!}} را توضیح دهید.
علاوه بر توضیح دستورها، به هر دو پرسش زیر در ویدیو پاسخ دهید:

\begin{enumerate}
  \item اگر دو نمونه از دستور \lr{\texttt{sleep}} را هم‌زمان اجرا
  کنیم، آیا این دو اجرا یک پردازه محسوب می‌شوند؟ کدام مشخصات آن‌ها می‌تواند
  یکسان باشد و کدام شناسه باید متفاوت باشد؟ اجرای آزمایش دوم لازم نیست.
  \item اگر علامت \lr{\texttt{\&}} را از دستور اول حذف کنیم، چه
  تغییری در امکان وارد کردن دستور بعدی در همان پایانه ایجاد می‌شود؟
  این تفاوت را با اجرای پیش‌زمینه و پس‌زمینه توضیح دهید.
\end{enumerate}

\paragraph{آزمون و ارزیابی}
در این سؤال، ساختار تحویل و وجود فایل‌های لازم به‌صورت خودکار بررسی
می‌شود. ارزیابی پاسخ‌ها بر اساس خروجی واقعی آزمایش، توضیح کوتاه گزارش
و پاسخ‌های ویدیو است. شناسه‌های عددی پردازه‌ها بین اجراها متفاوت‌اند
و لازم نیست با نمونه یا خروجی دانشجوی دیگری برابر باشند.

\paragraph{موارد لازم برای تحویل}
پاسخ این سؤال را در پوشهٔ \lr{\texttt{q3}}، مستقیماً در ریشهٔ آرشیو
یا فایل فشردهٔ تمرین و کنار پوشهٔ \lr{\texttt{q2}} قرار دهید.
ساختار این بخش به شکل زیر باشد:

\begin{latin}
\begin{verbatim}
q3/
  answers.md
  video.mp4
\end{verbatim}
\end{latin}

\begin{itemize}
  \item \lr{\texttt{answers.md}}: نسخهٔ هسته با خروجی
  \lr{\texttt{uname -r}}، دستورها و خروجی‌های آزمایش، پاسخ دو پرسش
  گزارش به فارسی و زمان شروع پاسخ دو پرسش ویدیو، مانند
  \lr{\texttt{01:10}}. همین فایل برای ثبت شواهد کافی است.
  \item \lr{\texttt{video.mp4}}: فایل ویدیوی پاسخ با قالب \lr{MP4}.
\end{itemize}

\textbf{نکته:} در صورت زیاد شدن حجم ویدیو می‌توانید با ابزار
\lr{ffmpeg} حجم آن را کاهش دهید و نسخهٔ کم‌حجم را قرار دهید.
